\subsection{The actual saturated resolved tail still has nonflat coefficient terms}

\begin{proposition}[A-L19: Proved; based saturation does not make these coefficient modules flat]
\label{AL19}
For every prime $p$, take the actual torsionfree
simplicial source $Y^{\mathrm{tf}}$ of A-P01, over
$R=\mathbb Z_{(p)}$ and $S=\Gamma_R(T)$, with
$b\geq3$ and $s=pb$. Let $\widetilde Y$ be its
entire iterated based PD saturation from A-SAT01
at cutoff $s+1$, with the specified actual map to
$G_R$. Then the weight-graded $S$-module
$\widetilde Y_2^0$ is not flat, even after forgetting
its grading. Thus this actual resolved-tail and
saturation construction does not supply the flat
coefficient terms used in A-L18.
\end{proposition}

\begin{proof}
We work in the specified simplex two and retain
the actual generators and free coefficient resolution
from (DP.3)--(DP.16). Put
\[
 a_0=p(p-1)+1,\qquad m_0=s+a_0,\qquad
 t=T,\quad c=(p-1)!.
 \tag{AL.97}
\]
In the rational ambient algebra give every resolved
generator from $W$ arity one and every old generator
arity zero. On zero-forms and their differentials,
this is an actual grading: old zero-forms belong to
$(G_R)_2^0$ and have old differentials, while
$d(W^0)\subset W^1$. The nonlinear twists in
(DP.16) occur on higher-form generators and do not
change this assertion about $d:Y_2^0\to Y_2^1$.
Products and all defined divided powers have their
indicated additive or multiplied arities.

The complete coefficient calculation in A-P01 gives
the following exact low-weight information:
there is no positive-arity rational zero-form of
arity one below weight $m_0$; in weight $m_0$ that
space is $\mathbb Q C$, with integral lattice $RC$.
The minimum weight in $W^1$ is $m_0-1$, where its
generator is the independent free generator $x$.
The actual formulas are
\[
 dC=pg+tx-pcy,\qquad
             t^{p-1}C=p(cK-H).
 \tag{AL.98}
\]
Here $g,y$ are independent positive-form generators
of weight $m_0$, and $K,H$ are the zero-forms
specified in (DP.11). These are identities in
$Y^{\mathrm{tf}}$, with no alteration of its twist.

We prove that the arity-one zero-form lattice in
weight $m_0$ remains $RC$ at every saturation stage.
The arity-zero zero-form algebra remains $(G_R)_2^0$.
Indeed its rational zero-forms have differentials
in the old $G_R$ one-form summand, and that summand
is a direct summand of the old form-one module.
The integral coefficient contraction of $G_R$
shows that any such rational zero-form with integral
derivative and integral vertex value was already
in $(G_R)_2^0$. Positive resolved arity cannot
contribute to this component. This proves the
arity-zero assertion at every successive stage.

There are no rational arity-one zero-forms below
$m_0$, at any stage, because all stages lie in the
same fixed rational ambient algebra. There are no
arity-two or higher zero-forms in weight $m_0$:
their minimum weight is $2m_0$. In weight $m_0$
a newly allowed arity-one zero-form can therefore
only be $rC$ for some $r\in\mathbb Q$.
Its derivative has $x$-coordinate $rt$ by (AL.98).
At the preceding stage this coordinate belongs
to the unchanged coefficient module
$(G_R)_2^{1,0}x$. No product of an arity-one
zero-form of weight at least $m_0$ with an old
positive one-form can contribute at weight $m_0$,
since those old forms have positive weight. No
higher-arity product contributes either.
Integrality of the derivative thus requires
$rt\in(G_R)_2^{1,0}$. Applying the actual vertex
map $(G_R)_2^{1,0}\to S^1=RT$ forces $r\in R$.
This proves the asserted lattice stability by
induction, including the PD closure after each
root-adjoining step and the filtered union.
The same argument proves arity homogeneity at each
stage: the derivative condition separates into
its arity components, and the vertex condition
restricts only arity zero.

Write $L=\widetilde Y_2^0$ and
$\overline L=L/pL$. The image $\overline C$ is
nonzero, because the arity-one weight-$m_0$ lattice
is exactly $RC$. It cannot belong to
$t\overline L$: taking the arity-one component
of weight $m_0$ would require an arity-one element
of weight $m_0-1$, whose group is zero. Yet (AL.98)
gives
\[
             t^{p-1}\overline C=0.
 \tag{AL.99}
\]

We finish with the precise flatness test. In
$D=S/pS=\Gamma_{\mathbb F_p}(T)$ one has
\[
             \ker(t^{p-1}:D\to D)=tD.
 \tag{AL.100}
\]
This can be verified directly on its divided-power
basis. Multiplication by $t$ sends $T^{[n]}$ to
$(n+1)T^{[n+1]}$, so its image is spanned by
$T^{[m]}$ with $m$ not divisible by $p$.
Multiplication by $t^{p-1}$ on $T^{[n]}$ has
coefficient $(n+1)\cdots(n+p-1)$, nonzero modulo
$p$ exactly when $n$ is divisible by $p$.
The images of the other basis elements vanish,
and distinct surviving basis elements have
distinct images. This proves (AL.100) for every
prime, including $p=2$.

Flatness of $L$ over $S$ would make
$\overline L=L\otimes_S D$ flat over $D$.
Tensoring the exact-at-the-middle sequence
\[
 D\xrightarrow{\ t\ }D
       \xrightarrow{\ t^{p-1}\ }D
 \tag{AL.101}
\]
with that flat module would give
$\ker(t^{p-1}:\overline L\to\overline L)
=t\overline L$. The explicitly verified element
$\overline C$ satisfies the left-hand condition
by (AL.99) and fails the right-hand condition.
This contradiction proves the stated nonflatness.
\end{proof}

The source in this proposition retains the actual
map, rational comparison, coefficient $H^0$ and
first skeleton established in A-P01 and A-SAT01.
The conclusion concerns its coefficient-module
flatness; a different replacement construction
with actual comparisons remains undecided.
